% ============================================================================
% Chapter 3: Types of Computers
% ============================================================================

\chapter{Types of Computers}
\label{ch:types}

\begin{objectives}
  \item Identify the four main types of personal computers
  \item Compare desktop, laptop, tablet, and smartphone features
  \item Understand which type of computer is best for different tasks
  \item Recognise that all these devices are computers
\end{objectives}

\bytetip[fun]{Did you know? The device in your pocket --- yes, your smartphone --- is a computer too! It just happens to fit in your jeans. Cool, right?}

% ---------------------------------------------------------------
\section{Different Types of Computers}
\label{sec:types-computers}
\index{types of computers}

Not all computers look the same. They come in many shapes and sizes, each designed for different purposes. Let's explore the four types you'll encounter most often:

% --- Figure 3.1: Types of Computers Visual Cards ---
\begin{figure}[H]
\centering
\begin{tabularx}{\textwidth}{@{}*{4}{>{\centering\arraybackslash}X}@{}}
  \begin{tcolorbox}[colback=PartColor1!8, colframe=PartColor1, width=3.4cm, height=4cm, valign=center, halign=center, arc=10pt]
    {\Huge\textcolor{PartColor1}{\faDesktop}}\\[6pt]
    \textbf{Desktop}\\[4pt]
    {\footnotesize Stays on a desk.\\Powerful.}
  \end{tcolorbox}
  &
  \begin{tcolorbox}[colback=PartColor3!8, colframe=PartColor3, width=3.4cm, height=4cm, valign=center, halign=center, arc=10pt]
    {\Huge\textcolor{PartColor3}{\faLaptop}}\\[6pt]
    \textbf{Laptop}\\[4pt]
    {\footnotesize Portable with\\a screen.}
  \end{tcolorbox}
  &
  \begin{tcolorbox}[colback=PartColor6!8, colframe=PartColor6, width=3.4cm, height=4cm, valign=center, halign=center, arc=10pt]
    {\Huge\textcolor{PartColor6}{\faBook}}\\[6pt]
    \textbf{Tablet}\\[4pt]
    {\footnotesize Touch screen.\\For reading.}
  \end{tcolorbox}
  &
  \begin{tcolorbox}[colback=PartColor7!8, colframe=PartColor7, width=3.4cm, height=4cm, valign=center, halign=center, arc=10pt]
    {\Huge\textcolor{PartColor7}{\faPhone}}\\[6pt]
    \textbf{Smartphone}\\[4pt]
    {\footnotesize Pocket-sized.\\Connected.}
  \end{tcolorbox}
\end{tabularx}
\caption{The four main types of personal computers}
\label{fig:types-cards}
\end{figure}

% ---------------------------------------------------------------
\section{Comparing Computer Types}
\label{sec:comparing-types}

Each type of computer has its own strengths and weaknesses. Let's compare them:

% --- Table 3.1: Computer Comparison Table ---
\begin{table}[H]
\centering
\caption{Comparing different types of computers}
\label{tab:compare-types}
\renewcommand{\arraystretch}{1.3}
\begin{tabularx}{\textwidth}{l X X X X}
\toprule
\rowcolor{HeaderRow}
\textcolor{white}{\textbf{Feature}} & 
\textcolor{white}{\textbf{Desktop}} & 
\textcolor{white}{\textbf{Laptop}} & 
\textcolor{white}{\textbf{Tablet}} & 
\textcolor{white}{\textbf{Smartphone}} \\
\midrule
\rowcolor{RowLight}
Portability & Low & Medium & High & Very High \\
\rowcolor{RowDark}
Screen Size & Large & Medium & Medium & Small \\
\rowcolor{RowLight}
Performance & High & Medium & Medium & Low--Medium \\
\rowcolor{RowDark}
Battery & \textcolor{WarnRed}{\ding{55}} No & \textcolor{LabGreen}{\ding{51}} Yes & \textcolor{LabGreen}{\ding{51}} Yes & \textcolor{LabGreen}{\ding{51}} Yes \\
\rowcolor{RowLight}
Best For & School lab, office & Travel, homework & Reading, browsing & Calls, messaging \\
\bottomrule
\end{tabularx}
\end{table}

\begin{notebox}
All four types --- desktop, laptop, tablet, and smartphone --- are \textbf{computers}. They all accept input, process data, store information, and produce output. They just do it in different shapes and sizes!
\end{notebox}

% ---------------------------------------------------------------
\section{Which Computer Should You Use?}
\label{sec:which-computer}

It depends on what you need to do! Here's a quick guide:

\begin{itemize}[leftmargin=*, label={\textcolor{ModuleBlue}{\faAngleRight}}, itemsep=4pt]
  \item \textbf{Writing a long essay?} Desktop or laptop --- a real keyboard makes typing faster.
  \item \textbf{Reading an e-book on the bus?} Tablet --- light, easy to hold, great screen.
  \item \textbf{Calling a friend?} Smartphone --- it's designed for communication.
  \item \textbf{Playing a high-end game?} Desktop --- most powerful option for gaming.
\end{itemize}

\begin{tryit}
Look around your classroom or home. Make a list of every computer-like device you can find. Don't forget things like smart TVs, gaming consoles, and smartwatches --- they're all computers too!
\end{tryit}

\bytetip[fun]{Smartphones are computers too! They just fit in your pocket. The phone you carry has more computing power than the computers used for the entire Apollo 11 Moon landing mission!}

\begin{funfact}
The first laptop was the \textbf{Osborne 1}, released in 1981. It weighed about \textbf{11 kilograms} (that's heavier than a bowling ball!) and had a tiny 5-inch screen. Today's laptops weigh less than 1.5 kg and have screens up to 17 inches. That's a LOT of progress!
\end{funfact}

% ---------------------------------------------------------------
% End-of-Chapter
% ---------------------------------------------------------------

\begin{reviewqs}
  \item Name four types of personal computers.
  \item Which type of computer is \emph{most portable}? Which is \emph{least portable}?
  \item Why would someone choose a desktop computer over a laptop?
  \item Is a smartphone a computer? Explain your answer.
  \item Give one example of a task best suited for each type of computer.
\end{reviewqs}

\begin{challenge}
Create a \textbf{poster} comparing the four types of computers. For each one, draw a simple picture, list 2 advantages and 1 disadvantage, and write one sentence about when you'd use it. Make it colourful!
\end{challenge}
